% !TEX program = xelatex
% =============================================================================
% cn-litigation · 刑事辩护词（一审）
%
% 刑事辩护词无官方格式，本样张结构对照刑辩专业所与司法行政机关公开发布的
% 辩护词通行体例：标题（辩护词＋案件引题）→ 称呼语（尊敬的审判长、审判员）
% → 受托与阅卷段 → 「一、」分点辩护意见（每点小标题即辩护要点结论）→
% 「综上所述」结语与量刑建议 → 此致法院 → 辩护人落款（律所＋律师＋汉字日期）。
%
% 样张：虚构交通肇事罪一审辩护（自首＋初犯＋赔偿谅解＋认罪认罚→建议缓刑），
% 所引法条均为现行有效条文。全部人名地名为占位。
%
% 编译：bash scripts/compile.sh templates/cn-litigation/defense-statement.tex --preview
% =============================================================================
\documentclass[zihao=-4]{ctexart}
\usepackage{litigation}

\begin{document}

\liTitle{辩\hspace{1\ccwd}护\hspace{1\ccwd}词}{——陈××涉嫌交通肇事罪一案}

\liVocative{尊敬的审判长、审判员：}

北京市示例律师事务所接受被告人陈××近亲属的委托，并经陈××本人同意，指
派王××律师担任其涉嫌交通肇事罪一案的一审辩护人。接受委托后，辩护人查阅
了本案全部卷宗材料，多次会见了被告人，并参加了今天的法庭调查，对本案事实
有了全面、客观的了解。

辩护人对公诉机关指控被告人陈××犯交通肇事罪的罪名不持异议。根据《中华人
民共和国刑事诉讼法》第三十七条关于辩护人职责的规定，辩护人依据事实和法律，
就被告人具有的法定、酌定从宽处罚情节发表如下辩护意见：

\section{被告人具有自首情节，依法可以从轻或者减轻处罚}

事故发生后，被告人立即停车、保护现场并拨打120急救电话和110报警电话，在现
场等候交通警察处理，到案后如实供述了事故经过。依据《中华人民共和国刑法》
第六十七条第一款及《最高人民法院关于处理自首和立功具体应用法律若干问题的
解释》第一条的规定，被告人的行为构成自首，依法可以从轻或者减轻处罚。

\section{被告人系初犯、偶犯，且属过失犯罪，主观恶性与人身危险性小}

被告人此前无任何违法犯罪记录，平时遵纪守法、表现良好，有村民委员会出具的
一贯表现证明在卷佐证。交通肇事罪系过失犯罪，被告人并非蔑视法秩序而积极追
求危害结果，其主观恶性和再犯可能性均明显小于故意犯罪。

\section{被告人已足额赔偿被害方全部损失并取得书面谅解}

案发后，被告人及其家属积极筹款，在保险理赔之外另行支付被害方各项损失合计
280,000元，被害方近亲属已出具《谅解书》，对被告人表示谅解并请求对其从宽处
罚。依据最高人民法院、最高人民检察院《关于常见犯罪的量刑指导意见》，积极
赔偿并取得谅解的，可以减少基准刑。

\section{被告人自愿认罪认罚，依法可以从宽处理}

被告人在审查起诉阶段自愿如实供述罪行，承认指控的犯罪事实，愿意接受处罚，
已在辩护人见证下签署《认罪认罚具结书》。依据《中华人民共和国刑事诉讼法》
第十五条的规定，可以依法从宽处理。

\vspace{0.3\baselineskip}
综上所述，被告人陈××犯罪情节较轻，具有自首、赔偿谅解、认罪认罚等多项从
宽处罚情节，有悔罪表现，没有再犯罪的危险，宣告缓刑对所居住社区没有重大不
良影响。辩护人恳请合议庭充分考虑上述情节，依据《中华人民共和国刑法》第七
十二条之规定，对被告人陈××从轻处罚并适用缓刑，给其改过自新、回报社会的
机会。

以上辩护意见，恳请合议庭评议时充分考虑并予采纳。

\liSubmitTo{示例市示例区人民法院}

\liSign{辩护人：北京市示例律师事务所}
\liSign{王××\hspace{1\ccwd}律师}
\liSign{二〇二六年七月十八日}

\end{document}
